\section{A100的急变、软性与低能响应}\label{sec:4-3}

如果说 $A\approx 60$ 与 $80$ 区域展现了壳模型次闭壳破缺与多极小构型混合的生动图像，那么将视线转向 $A\approx 100$ 核区（主要涵盖 $\mathrm{Zr}$、$Z=42\,\mathrm{Mo}$、$Z=44\,\mathrm{Ru}$ 与 $Z=46\,\mathrm{Pd}$ 同位素），核结构演化则呈现出更为剧烈、尺度交织的多重动力学特征。长期以来，该区域在核物理文献中占据着双重经典地位：一方面，锆（$\mathrm{Zr}$）同位素链在中子数 $N=56\text{--}60$ 处被誉为核多体系统中常被讨论的急剧``一阶形状量子相变''（Quantum Phase Transition, QPT）范例；另一方面，邻近的 $\mathrm{Mo}$ 与 $\mathrm{Ru}$ 同位素则常被教科书作为相互作用玻色子模型（IBM）中 $O(6)$ 动力学对称性（即 $\gamma$ 软核）的标准坐标。

然而，将整个 $A\approx 100$ 核区均质化为单一 $O(6)$ 对称性极限，是一种忽略核微观异质性的严重误区。本节将基于 Togashi 等（2016）\cite{doctoral-r1-g28-togashi-2016-zr-mcsm-qpt}的大规模蒙特卡罗壳模型（MCSM）计算、Singh 等（2018）\cite{doctoral-r1-g28-singh-2018-98zr-lifetimes}的寿命实验原件、Chen \& Singh（2020）\cite{chen-singh-2020-nds-98mo}与 Thomas 等（2016）\cite{thomas-2016-96-98mo-coexistence}关于 $^{98}\mathrm{Mo}$ 的推荐能级与反应谱学、Karmakar 等（2026）\cite{karmakar-2026-98mo-octupole}的高自旋出版全文，以及 Saracino 等（2026）\cite{saracino-2026-102ru-scissors-octupole-resonances}关于 $^{102}\mathrm{Ru}$ 的核共振荧光（NRF）测量，系统阐明质量区内部由质子与中子亚壳填充驱动的非均质势能面演化，并严格解耦低能四极 $\gamma$ 软性、等矢 $M1$ 剪刀模与反射不对称八极共振之间的物理尺度界限。

\subsection{Zr同位素链的一阶形状量子相变与Type-II壳演化}\label{sec:4-3-1}

在偶偶核 $\mathrm{Zr}$ 同位素链（$Z=40$）中，基态电四极形变的演化展现出明显的急变。在 $N=50$（$^{90}\mathrm{Zr}$）至 $N=56$（$^{96}\mathrm{Zr}$）区间，基态呈现高度刚性的球形闭壳特征，第一激发 $2^+_1$ 态能量居高不下（例如 $^{96}\mathrm{Zr}$ 中 $E(2^+_1) = 1750\,\mathrm{keV}$）；中子数仅增加 4 颗至 $N=60$（$^{100}\mathrm{Zr}$）时，第一 $2^+_1$ 态能量急剧坠落至仅 $213\,\mathrm{keV}$，同时约化跃迁几率 $B(E2; 0^+_1 \to 2^+_1)$ 暴增数十倍，在指定MCSM分析中对应较大形变成分的出现，而几何指认仍需结合该模型与电磁观测。

这一急剧相变的微观物理本质曾困扰学界数十年。宏观几何模型或唯象玻色子模型仅能通过人为调节势能面参数强制拟合这一突变，而无法给出微观因果机制。Togashi 等（2016）\cite{doctoral-r1-g28-togashi-2016-zr-mcsm-qpt}利用大规模蒙特卡罗壳模型（MCSM），在包含质子与中子完整主要价空间的超大组态空间内，成功揭示了驱动这一相变的微观引擎——第二类壳演化（Type-II shell evolution）：
\begin{enumerate}
\item 质子闭壳弱化与跨壳激发：在球形基态下，质子闭合于 $Z=40$ 次闭壳（主要填满 $1f_{5/2}, 2p_{3/2}, 2p_{1/2}$ 轨道），此时向上跃迁至质子 $1g_{9/2}$ 独特轨道的激发能极高；
\item 张量力与自旋-同位旋单极吸引：随着中子数填充至 $N=60$，中子开始占据 $1g_{7/2}$ 与 $1h_{11/2}$ 轨道。质子 $1g_{9/2}$ 轨道与中子 $1g_{7/2}$ 轨道之间存在极其强烈的单极吸引相互作用（核子间反向自旋张量力分量与同位旋矢量交换力，见第~2.2~节式~\eqref{eq:2-2-1}）；
\item 自强化形变正反馈：质子向 $1g_{9/2}$ 轨道的少量热激发或组态混合，通过强烈的单极积分使得中子 $1g_{7/2}$ 与 $1h_{11/2}$ 的有效单粒子能量（ESPE）急剧下移，反过来进一步大幅降低了质子跨壳激发的能垒。
\end{enumerate}
这种占据诱导能隙坍塌的非线性正反馈，使得 $^{100}\mathrm{Zr}$ 的球形构型瞬间转变为极深的大长椭极小，从而微观实现了多体系统中的一阶形状量子相变。

紧邻相变临界点的 $^{98}\mathrm{Zr}$（$N=58$）展现了相变前夕的微观构型共存特征。核对 Singh 等（2018）\cite{doctoral-r1-g28-singh-2018-98zr-lifetimes}反冲距离柱塞法（RDDS）原件 PDF 第 3 页 Table I：实验测得第一激发 $2^+_1$ 态寿命为 $\tau = 3.8(8)\,\mathrm{ps}$，其向基态跃迁强度仅为 $B(E2; 2^+_1 \to 0^+_1) = 2.9(6)\,\mathrm{W.u.}$，直接证实了基态仍受闭壳约束、具有极低的四极集体性；然而同一 $2^+_1$ 态向激发态 $0^+_2$（853 keV）的退激分支表现出显著增强的集体性（$B(E2) = 28.3(60)\,\mathrm{W.u.}$），且随自旋升高，晕态带内跃迁强度急剧上升至 $B(E2; 6^+_1 \to 4^+_1) = 103(36)\,\mathrm{W.u.}$。这清楚地证明了在相变临界点附近，近球形构型与大形变长椭构型共存，且两者之间存在阻碍基态带内直接强烈跃迁的明显势能阻尼。

\subsection{质量区内N与Z的非均质性：从Zr到Mo、Ru的势能演化与IBM O(6)极限适用边界}\label{sec:4-3-2}

Zr 同位素链所展现的极端陡峭一阶相变，常常给读者造成一种错觉，仿佛整个 $A\approx 100$ 核区都在经历相同模式的形状相变。然而，深入考察质子数 $Z$ 的微小改变，便会发现质量区内部强烈的物理非均质性（heterogeneity）：
\begin{itemize}
\item \textbf{$Z=40$（$\mathrm{Zr}$ 链）}：质子严格受制于 $Z=40$ 次壳隙，核芯极化受阻，表现为球形与超强形变长椭球之间的极端对立；
\item \textbf{$Z=42$（$\mathrm{Mo}$ 链）}：仅增加 2 个质子进入 $1g_{9/2}$ 轨道，费米面附近的质子简并度骤增，$Z=40$ 亚壳隙被彻底冲垮，势能面对四极形变呈现出极大的平坦性与宽软性；
\item \textbf{$Z=44$（$\mathrm{Ru}$ 链）与 $Z=46$（$\mathrm{Pd}$ 链）}：质子进一步向 $1g_{9/2}$ 中部填充，质子粒子与中子空穴的横向竞争使得势能面在 $\gamma$ 方向展现出稳定的三轴凹陷或平坦的 $\gamma$ 不稳定性。
\end{itemize}

相互作用玻色子模型（IBM-1）在描述 $A\approx 100$ 核区时，通常采用六维正交群 $O(6)$ 动力学对称性链。在理想 $O(6)$ 极限下，其相干态经典极限对应的势能面关于三轴角 $\gamma$ 完全无规平坦（即理想 Wilets-Jean 软核）。

然而，将理想 $O(6)$ 对称性生搬硬套到整个 $A\approx 100$ 区域面临严重的微观失效。微观自洽平均场与协变密度泛函（CDFT）计算表明，随着中子数向 $N=60$ 靠近，无论是 $\mathrm{Mo}$ 还是 $\mathrm{Ru}$，其势能面均显现出局域微观壳效应所雕刻的微弱三轴极小，且在 $\beta$ 方向的刚度远大于理想玻色子势。更重要的是，唯象 IBM-1 忽略了质子与中子玻色子的独立自由度，无法容纳等矢剪刀振动模式与微观准粒子对拆裂，因而绝不能将 $A\approx 100$ 概括为均质化的理想 $O(6)$ 世界。

\subsection{98Mo的低能结构与形状共存：$0^+_2$态性质、E0单极跃迁与八极关联初探}\label{sec:4-3-3}

作为非均质演化的核心坐标，偶偶核 $^{98}\mathrm{Mo}$（$Z=42, N=56$）的低激发结构集中体现了四极形状共存与大振幅软性的复杂交织。

查阅 Chen \& Singh（2020）\cite{chen-singh-2020-nds-98mo}的最新 $A=98$ 推荐能级原件（NDS 164, 1，PDF 第 2 页）：
\begin{itemize}
\item 第一激发态为 $2^+_1$ 态，准确能量为 $E = 787.384(17)\,\mathrm{keV}$，其实测半衰期为 $T_{1/2} = 3.47(7)\,\mathrm{ps}$（对应平均寿命 $\tau = 5.01(10)\,\mathrm{ps}$）；
\item 低能 $0^+_2$ 态，准确能量位于 $E = 734.75(4)\,\mathrm{keV}$（其能量略低于第一 $2^+_1$ 态）；该态半衰期测定为纳秒量级的 $T_{1/2} = 21.8(9)\,\mathrm{ns}$（派生平均寿命 $\tau = T_{1/2}/\ln 2 = 31.45(130)\,\mathrm{ns}$）；
\item 在退激通道上，由于 $734.75\,\mathrm{keV} < 787.384\,\mathrm{keV}$，能级运动学严格禁止 $0^+_2 \to 2^+_1$ 的向下放 $\gamma$ 衰变；相反，Thomas 等（2016）\cite{thomas-2016-96-98mo-coexistence}的Table 2列入了 $2^+_1\to0^+_2$ 的约52.6 keV支路，但以双星号标为NDS采用量，且强度栏的横线表示该次实验未观测到该线，而 $0^+_2$ 态自身则通过内转换电子向基态发射单极 $E0$ 跃迁。
\end{itemize}

文献以构型混合描述这些低能态时，讨论近球形成分与较大形变成分的相互作用。这是对能级、分支和矩阵元的模型解释；仅凭低能 $0_2^+$ 与相对跃迁强度，不能唯一确定两个极小之间的势垒或完整 $\gamma$ 势面。与Ni案例比较时，应先对齐采用的混合假设和观测量。

更为前沿的高自旋探索来自于近期发表的完整论文 Karmakar 等（2026）\cite{karmakar-2026-98mo-octupole}。核对该 8 页出版论文原件（Nucl. Phys. A 1066, 123265）：
\begin{itemize}
\item 原件 PDF 第 4 页 Table 1 汇报了两条新发现的退激跃迁线：\allowbreak{}$1833.3(5)\,\mathrm{keV}$ 跃迁（初态为 $2620.4\,\mathrm{keV}\,[5^-]$，终态为 $787.4\,\mathrm{keV}\,[2^+_1]$）\allowbreak{}以及 $1585.8(6)\,\mathrm{keV}$ 跃迁（初态为 $3095.4\,\mathrm{keV}\,[7^-]$，终态为 $1509.6\,\mathrm{keV}\,[4^+_1]$）。因初末态自旋宇称差，表中将多极性标注为最低允许的 $(E3)$，但必须注意，原表中这两条线的角分布\allowbreak{} $R_{\mathrm{DCO}}$ 与偏振不对称度 $\Delta_{\mathrm{iPDCO}}$ 数据均为空白，实验并未独立测定其多极性；
\item 原件 PDF 第 6 页 Table 2 给出了电八极与电四极跃迁率比值：对于 $5^-$ 态，测得分支比导出 $B(E3)/B(E2) = 66.05 \pm 27.01\,\mathrm{fm}^2$；对于 $7^-$ 态，导出 $B(E3)/B(E2) = 153.81 \pm 33.31\,\mathrm{fm}^2$；
\item 原文对比指出，上述 $B(E3)/B(E2)$ 比值相较于 $^{100}\mathrm{Ru}$ 约高出 4 倍；然而必须明确界定其物理边界：该实验并未测定 $5^-$ 与 $7^-$ 能级的绝对寿命，所给出的仅是分支比导出的比值，\textbf{绝不能将比值高出 4 倍误读为绝对 $B(E3)$ 强度高出 4 倍}。
\end{itemize}
因此，在 $^{98}\mathrm{Mo}$ 中确立的是低能四极形状共存以及高自旋下增强的负宇称八极关联线索，绝不能跳步宣布已经证实了高自旋内禀稳定八极形变。

\subsection{102Ru的核共振荧光测量：M1剪刀模与E1八极共振的共存}\label{sec:4-3-4}

进入钌同位素，偶偶核 $^{102}\mathrm{Ru}$（$Z=44, N=58$）展现出更为丰富的多极电磁激发现象。Saracino 等（2026）\cite{saracino-2026-102ru-scissors-octupole-resonances}利用美国北卡罗来纳高等 $\gamma$ 射线源（$\mathrm{HI}\vec{\gamma}\mathrm{S}$）装置，通过准单能、高通量线偏振光子束，对 $^{102}\mathrm{Ru}$ 开展了精密的核共振荧光（Nuclear Resonance Fluorescence, NRF）光核散射实验。

NRF 实验通过测量光子弹性散射角分布与方位角线偏振不对称度，能够直接提取激发态的宇称与绝对衰变线宽 $\Gamma_0$ 及向基态跃迁的约化几率。核对 Saracino 等（2026）原件：
\begin{enumerate}
\item \textbf{低能偶极 $E1$ 跃迁分布}：原件 PDF 第 3 页明确汇报，在 $E_\gamma = 2.95\,\mathrm{MeV}$ 处观察到一条非常显著的单孤立 $E1$ 跃迁，其绝对跃迁强度测定为 $B(E1)\uparrow = 0.016\,e^2\mathrm{fm}^2$；同时在更高激发能 $E_\gamma \approx 4.5\text{--}5.5\,\mathrm{MeV}$ 区间（中心位于 $4.9\,\mathrm{MeV}$ 附近），观察到共振峰状聚集的电偶极吸收，其求和跃迁强度达到 $\sum B(E1)\uparrow \approx 0.1\,e^2\mathrm{fm}^2$；
\item \textbf{等矢磁偶极剪刀模（Scissors Mode）}：原件 PDF 第 5 页指出，实验测得的总观测磁偶极强度达到 $\sum B(M1)\uparrow = 1.7\,\mu_N^2$，这与形变度 $\beta = 0.20$ 的求和规则预期高度一致；实验提取的光谱密度分布 $S(M1; E)\uparrow$ 展现出中心位于 $3.6\,\mathrm{MeV}$ 附近的宽共振轮廓，符合剪刀模式的标准系统发育；而在 $2\text{--}4\,\mathrm{MeV}$ 的低能局域区间，观测到的局部剪刀模强度约为 $1.0\,\mu_N^2$；
\item \textbf{八极共振的微观模型解释}：作者对 $2.95\,\mathrm{MeV}$ 孤立 $E1$ 态提出可能的四极—八极双声子解释，并比较声子和能量 $E(2_1^+)+E(3_1^-)=2518\,\mathrm{keV}$；约 $4.9\text{--}5.0\,\mathrm{MeV}$ 的聚集则结合其TAC-QRPA计算解释为四极形变下的偶极—八极混合及低能八极共振。两项解释的模型范围和确定程度不同，均须与直接报告的偶极强度分开。
\end{enumerate}
Saracino 等（2026）的测量确凿证明：在一个中等质量过渡核中，磁偶极剪刀振荡与电偶极/八极共振响应完全可以在相邻的能量窗口内稳定共生。

\subsection{不同物理尺度响应的严格解耦：低能γ软性、剪刀模与八极激发不可混同}\label{sec:4-3-5}

面对 $^{102}\mathrm{Ru}$ 与 $^{98}\mathrm{Mo}$ 呈现的多重模式，必须确立最为严格的认识论解耦红线。文献中常有研究者试图用单一的几何形变参量（如引入某种静态八极三轴形变 $\beta_3, \gamma$），将低能能谱的 $\gamma$ 软性、剪刀模以及偶极跃迁打包解释为同一种序参量的破缺。这种混淆抹杀了三个现象在物理本质上的根本分野：
\begin{itemize}
\item \textbf{低能 $\gamma$ 软性（Quadrupole $\gamma$-Softness）}：
  讨论的是四极形状自由度沿 $\gamma$ 方向的势能、动能与量子分布；不能由固定能窗定义软性。其波函数与电四极算符共同决定相关能谱、矩阵元和分支；
\item \textbf{剪刀振动模式（Scissors Mode）}：
  物理尺度位于 $2.5\text{--}4.0\,\mathrm{MeV}$，属于\textbf{等矢一阶磁偶极内部角向自由度}。它依赖于质子流体与中子流体之间的相对角动量恢复力，要求核必须具有非零内禀四极形变，但并不要求静态三轴性或宇称破缺，其探针严格为绝对磁偶极强度 $B(M1)$；
\item \textbf{八极共振与负宇称态（Octupole Resonances）}：
  物理尺度横跨 $2.0\text{--}5.5\,\mathrm{MeV}$，属于\textbf{三阶八极表面电荷反射不对称自由度}。在 $^{102}\mathrm{Ru}$ 中观察到的主要是低能 $2.95\,\mathrm{MeV}$ 处的双声子振动及 $4.9\,\mathrm{MeV}$ 处的低能巨共振低翼分布，其解释是作者结合多体理论给出的动力学模型结论，并非直接测得静态内禀八极形变 $E3$ 矩阵元。
\end{itemize}

三者涉及的相互作用通道（四极配对力、等矢自旋-同位旋偶极力、八极粒子-空穴力）、微观算符对称性以及激发的物理自由度完全独立，绝非同一种序参量。

更重要的是，必须划定实验探针的推论红线：从 $0^+$ 基态出发的光核吸收反应虽然绝大部分截面集中于 $E1$ 与 $M1$ 偶极跃迁（布居 $J^\pi = 1^\pm$ 态），但在电四极光子吸收下同样能够布居少量的 $J^\pi = 2^+$ 态；不存在所谓 NRF 只能唯一激发自旋 1 的绝对物理定理。然而，其实测态的自旋严格受限于低多极光子吸收。因此，在 $^{102}\mathrm{Ru}$ 中观察到 $1^+$ 与 $1^-$ 态的丰富结构，仅仅提供了低自旋偶极强度的能谱分布证据；绝不能跳步将这种低能偶极共振的共存，直接当作高自旋手征四重带或转动摆动模式的认证证明。检验高自旋集体模式还需相关自旋窗口内的能带连通、跃迁性质、寿命及模型联合预言，不能由低自旋的单一偶极观测外推。
